Reflexivní portfolia kladou důraz na výběr 3--5 reprezentativních prací doplněných o metareflexi a dokumentaci postupu; cílem je hodnotit schopnost vysvětlit postup, rozhodnutí a případné použití AI.

Praktický postup: 
\begin{enumerate}
\item Návrh: povinné komponenty (3--5 artefaktů), dokumentace postupu (drafty, mezikroky, screenshoty, časová razítka) a metareflexe.
\item Technické odevzdání: sběr v prostředí zachovávajícím mezikroky (např. OneNote) nebo upload do LMS; OneNote usnadňuje převod rukopisu a sběr postupů \cite{kb_ai_101_tipu_pdf_04e4a115_page_8_1}.
\item Instrukce a kontrola: přiložte šablonu metareflexe, pravidla pro deklaraci AI (typ, prompt, úpravy) a použijte spot‑check pro ověření autenticity.
\end{enumerate}

Minimum portfolia: 3--5 prací; u každé první draft, revidovaná verze/poznámky a chronologický záznam; povinná metareflexe a explicitní deklarace AI.

Rubrika (procesně orientovaná): dokumentace postupu 35\,\% \cite{kb_ai_101_tipu_pdf_04e4a115_page_8_1}, kvalita metareflexe 35\,\%, učební pokrok 20\,\%, formální náležitosti 10\,\%.

Šablona metareflexe: 
\begin{enumerate}
\item Co jsem dělal(a) a proč byl artefakt vybrán? (1--2 odstavce) \item Které AI nástroje jsem použil(a) a jak (prompt/funkce)? Jak jsem ověřil(a) výstupy?
\item Co jsem se naučil(a) mezi verzemi — hlavní zlepšení a nezodpovězené otázky. \item Krátký plán dalšího kroku (1--2 věty).
\end{enumerate}

Poznámka: požadujte artefakty dokazující proces (drafty, screenshoty, časová razítka); OneNote a LMS usnadní sběr a ověření \cite{kb_ai_101_tipu_pdf_04e4a115_page_8_1}.